\documentclass[11pt]{article}
\usepackage[margin=1in]{geometry}
\usepackage{amsmath,amssymb,amsthm,hyperref}
\newtheorem{theorem}{Theorem}
\title{A complete market has one martingale vertex, not zero}
\author{ziangni-sys}
\date{Conjecture 00000008251}
\begin{document}
\sloppy
\maketitle
\noindent The conjecture asserts that the number of extreme points of the
martingale-measure set of a finite-state market equals the number of
states minus the rank of the assets. A complete two-state market already
contradicts this formula.

\section*{An explicit complete, arbitrage-free market}
Take $\Omega=\{0,1\}$ with all subsets measurable and physical probability
$P=\tfrac12\delta_0+\tfrac12\delta_1$. There are two dates, $0$ and $1$,
with trivial initial information and full terminal information. Use
discounted prices and two Arrow securities:
\[
 S_0^0=S_0^1=\tfrac12,\qquad S_1^i(x)=\mathbf1_{\{i\}}(x).
\]
The terminal payoff matrix, with states indexing rows and securities
indexing columns, is
\[
 D=\begin{pmatrix}1&0\\0&1\end{pmatrix},\qquad
 \operatorname{rank}D=2=|\Omega|.
\]
The portfolio $(h_0,h_1)$ pays $h_x$ in state $x$ and costs
$(h_0+h_1)/2$, so every contingent claim is replicated. The portfolio
$(1,1)$ costs and pays one, supplying a constant numeraire.
Nonnegative payoffs and nonpositive cost force both coordinates to be
zero, proving absence of arbitrage.

\medskip
\noindent\textbf{All martingale measures.}
For a probability measure $Q$ on this actual measurable space, the
one-period discounted martingale equations are exactly
\[
 \mathbb E_Q[S_1^i]=S_0^i\quad(i=0,1).
\]
Conditional expectation at the trivial initial sigma-algebra is ordinary
expectation. Integrating the singleton indicators gives
\[
 Q(\{0\})=Q(\{1\})=\tfrac12.
\]
Singleton masses determine every measure on this finite space. Therefore
$Q=P$ is the only martingale probability. It is equivalent to $P$, since
both masses are positive. Thus the set of all martingale probabilities
and the set of equivalent martingale probabilities are both the singleton
$\{P\}$.

\begin{theorem}
The conjectured extreme-point formula fails for this market.
\end{theorem}
\begin{proof}
The unique point of a singleton convex set is extreme: any convex
decomposition into points of that set has both terms equal to that point.
Consequently either convention for martingale measures gives exactly one
extreme point. The proposed formula instead gives $2-2=0$.
\end{proof}

\newpage
\raggedright
\section*{Scope}
Both originals state the count formula for a finite-state market without
requiring incompleteness. The separate assertion about replication in
incomplete markets is a different conjunct. This counterexample is a
nonempty, complete, arbitrage-free market with a strictly positive
martingale probability, so neither nonexistence nor boundary probability
conventions explain the failure. The Arrow-security payoff matrix has
rank two directly; the argument does not depend on whether a separately
named bond is counted among the assets. The number of extreme points is
not the affine dimension of the feasible set: the singleton has dimension
zero but one extreme point.

\section*{Actual probability and financial objects in Lean}
The state type is \texttt{Fin 2} with its actual discrete Borel measurable
structure. \texttt{payoff}, \texttt{price}, \texttt{terminal} and
\texttt{cost} define the two securities and actual portfolio sums.
\texttt{terminal\_eq} proves that a portfolio pays its corresponding
coordinate in each state. \texttt{complete\_market},
\texttt{bond\_replication} and \texttt{no\_arbitrage} establish the
financial properties above. \texttt{payoffMatrix} is the actual matrix
of these payoffs; it is proved to equal the identity, and its actual
\texttt{Matrix.rank}, defined from the range of its linear map, is two.

\texttt{rawReference} is the sum of two genuinely scaled Dirac measures.
Its total mass is proved to be one, giving the bundled
\texttt{ProbabilityMeasure} called \texttt{reference}.
\texttt{MartingaleMeasure} quantifies over both actual Bochner integrals
of the securities and equates them to the initial prices.
\texttt{expectation\_eq\_weight} proves the singleton-expectation formula
from indicator integration. The definition of martingale measure thus
uses the complete finite-market equations, not asserted scalar weights.

\texttt{weight\_injective} uses the measure-extensionality theorem on
countable spaces: agreement of all real singleton masses implies equality
of the actual probability measures. \texttt{martingale\_unique} proves
an if-and-only-if characterization of every martingale probability.
\texttt{equivalent\_martingale\_unique} also imposes both genuine
absolute-continuity conditions, and proves the same characterization.

The finite mass coordinates place the full sets in the real vector space
$\mathbb R^\Omega$. \texttt{martingaleWeights} and
\texttt{equivalentWeights} are defined as images of all the actual
probability measures satisfying the corresponding conditions. Both sets
are proved to be the singleton uniform vector. The final count uses
Mathlib's actual \texttt{extremePoints} and \texttt{Set.ncard}.
The final theorem compares this computed count with the actual state
cardinality minus the actual payoff-matrix rank. No probability, rank,
feasible-set size or extreme-point count is assumed.

\section*{Reproduction}
The project pins Lean 4.19.0 and Mathlib commit
\begin{center}\small\texttt{c44e0c8ee63ca166450922a373c7409c5d26b00b}.\end{center}
From \texttt{lean/}, run \texttt{lake update}, \texttt{lake exe cache get},
\texttt{lake build}, and
\texttt{lake env lean Main.lean -DwarningAsError=true}.
The printed theorem dependency audits contain only standard logical
axioms; no admitted statements or auxiliary computations are used.
\end{document}
